\documentclass[12pt,a4paper]{article}

\usepackage{fontspec}
\usepackage{polyglossia}
\usepackage{geometry}
\usepackage{array}
\usepackage{tabularx}
\usepackage{longtable}
\usepackage{enumitem}
\usepackage{ragged2e}

\setmainlanguage{russian}
\setmainfont{Times New Roman}

\geometry{left=30mm,right=15mm,top=20mm,bottom=20mm}
\pagestyle{empty}
\setlength{\parindent}{0pt}
\setlength{\parskip}{0pt}
\renewcommand{\arraystretch}{1.18}
\setlist{nosep,leftmargin=*}

\newcommand{\field}[1]{\underline{\makebox[\linewidth][c]{#1}}}
\newcommand{\sigline}[1]{\underline{\hspace{42mm}}\,/ #1 /}

\begin{document}

\begin{center}
\footnotesize
ФЕДЕРАЛЬНОЕ ГОСУДАРСТВЕННОЕ АВТОНОМНОЕ ОБРАЗОВАТЕЛЬНОЕ УЧРЕЖДЕНИЕ\\
ВЫСШЕГО ОБРАЗОВАНИЯ\\
\textbf{«САНКТ-ПЕТЕРБУРГСКИЙ ПОЛИТЕХНИЧЕСКИЙ УНИВЕРСИТЕТ ПЕТРА ВЕЛИКОГО»}
\end{center}

\vspace{0.5em}

\begin{center}
\begin{tabularx}{\textwidth}{@{}>{\centering\arraybackslash}X@{}}
\hline
Высшая школа прикладной математики и вычислительной физики\\
\hline
\footnotesize (наименование учебного подразделения)\\
\end{tabularx}
\end{center}

\vspace{1.0em}

\begin{center}
\textbf{ИНДИВИДУАЛЬНЫЙ ПЛАН (ЗАДАНИЕ И ГРАФИК)}\\
\textbf{ПРОВЕДЕНИЯ ПРАКТИКИ}
\end{center}

\vspace{0.8em}

\field{Уртемеев Сергей Александрович}

\vspace{0.8em}

\begin{tabularx}{\textwidth}{|X|}
\hline
Тема практики: Расчет загрузки маршрутов общественного транспорта по расписанию\\
\hline
Направление подготовки (код/наименование) 01.03.02 «Прикладная математика и информатика»\\
\hline
Профиль (код/наименование) 01.03.02\_01 «Математическое моделирование и искусственный интеллект»\\
\hline
Вид практики Производственная\\
\hline
Тип практики Преддипломная\\
\hline
Место прохождения практики СПбПУ, Высшая школа прикладной математики и вычислительной физики; Политехническая ул., 29Б, Санкт-Петербург\\
\hline
\end{tabularx}

\vspace{0.8em}

\begin{tabularx}{\textwidth}{|>{\centering\arraybackslash}X|}
\hline
Руководитель практики от ФГАОУ ВО «СПбПУ»:\\
\hline
Курц В.\,В., канд. физ.-мат. наук, доцент\\
\hline
\end{tabularx}

\begin{center}
\footnotesize (Ф.И.О., уч. степень, должность)
\end{center}

\begin{tabularx}{\textwidth}{|X|}
\hline
Руководитель практики: Курц В.\,В., канд. физ.-мат. наук, доцент\\
\hline
\end{tabularx}

\begin{center}
\footnotesize (Ф.И.О., должность)
\end{center}

\vspace{0.8em}

\begin{center}
\textbf{Рабочий график проведения практики}
\end{center}

Сроки практики: 02.02.2026 по 07.05.2026

\vspace{0.6em}

\small
\begin{longtable}{|>{\centering\arraybackslash}p{0.07\textwidth}|p{0.18\textwidth}|p{0.45\textwidth}|p{0.15\textwidth}|}
\hline
№ п/п & Этапы (периоды) практики & Вид работ & Сроки прохождения этапа\\
\hline
1 & Организацион\-ный этап &
Установочная консультация; уточнение целей преддипломной практики; согласование индивидуального плана; определение состава материалов, связанных с ВКР. &
02.02.2026\\
\hline
2 & Основной этап &
2.1. Анализ предметной области: расчет загрузки маршрутов общественного транспорта по расписанию, OD-спрос, пассажирские соединения, ограничения пересадок. &
03.02.2026--20.02.2026\\
\cline{3-4}
 & &
2.2. Формализация модели: транспортное предложение, сегменты маршрутов и соединений, временные интервалы спроса, метрики соединения, обобщенное сопротивление. &
23.02.2026--13.03.2026\\
\cline{3-4}
 & &
2.3. Изучение и доработка программной реализации: загрузка входных таблиц, предварительная обработка сети, поиск соединений методом ветвей и границ, отбор альтернатив и распределение спроса. &
16.03.2026--17.04.2026\\
\cline{3-4}
 & &
2.4. Проверочные расчеты: подготовка синтетического набора данных, контроль сохранения спроса, нормировки вероятностей, согласованности нагрузок и выходных таблиц. &
20.04.2026--30.04.2026\\
\hline
3 & Заключительный этап &
Подготовка отчетных материалов по практике; оформление результатов для использования в ВКР; представление отчета по практике. &
01.05.2026--07.05.2026\\
\hline
\end{longtable}

\normalsize
\vspace{1.0em}

\begin{tabularx}{\textwidth}{@{}Xr@{}}
Обучающийся & \sigline{Уртемеев С.\,А.}\\[1.2em]
Руководитель практики от ФГАОУ ВО «СПбПУ» & \sigline{Курц В.\,В.}\\[1.2em]
Согласовано: руководитель практики & \sigline{Курц В.\,В.}\\
\end{tabularx}

\end{document}
